\documentclass[12pt]{article}
\usepackage[utf8]{inputenc}
\usepackage[T1]{fontenc}
\usepackage{hyperref}
\title{Rawa Language Resources for K-12 Curriculum Development}
\author{Compiled by Claude Sonnet 5.5 (claude-sonnet-5-5) for Andr\'as Kornai}
\date{}
\begin{document}
\maketitle

\section{Introduction}
Rawa (ISO 639-3 \texttt{rwo}; also Erawa, Erewa, Raua, Karo) is a Finisterre language of Papua New Guinea. It has two dialects, Rawa and Karo, on opposite sides of the Finisterre Range. Wikipedia places it in Morobe Province, while a SIL numerals dataset places about 11,000 speakers in the Rai Coast and Upper Ramu districts of Madang Province. The sources disagree on the province.

\section{Primary Educational Resources}
I found no primers, textbooks or school programmes in Rawa.

\section{Linguistic and Reference Materials}
\begin{itemize}
\item Wikipedia, Rawa language (phonology, dialects): \url{https://en.wikipedia.org/wiki/Rawa_language}
\item Numerals dataset entry for Rawa (Max Planck Institute for Evolutionary Anthropology): \url{https://lingweb.eva.mpg.de/channumerals/Rawa.htm}
\end{itemize}
The main reference grammar is Toland and Toland (1991), \emph{Reference grammar of the Karo/Rawa language} (SIL, Ukarumpa). I found it cited but saw no online copy.

\section{Religious and Community Resources}
\begin{itemize}
\item \emph{Anutro Mande Keta}, Rawa-dialect New Testament (1992), eBible page: \url{https://ebible.org/study/content/texts/rwo-rawa/about.html}
\item \emph{Anutro Mande Keta}, copyright and licence page (CC BY-ND 4.0): \url{https://ebible.org/rwo-rawa/copyright.htm}
\item \emph{Anutlo Sumange Mande Keda}, Karo-dialect New Testament (2001), eBible text: \url{https://ebible.org/rwo-karo/FRT01.htm}
\item \emph{Anutlo Sumange Mande Keda}, eBible page: \url{https://ebible.org/study/content/texts/rwo-karo/about.html}
\item Digital Bible Society entry: \url{https://dbs.org/bibles/RWOWBT}
\end{itemize}

\section{Digital and Audio Resources}
No audio resources were found. The online New Testament texts are the only digital texts I found.

\section{Curriculum Development Context}
The two New Testaments in the Rawa and Karo dialects are the largest known bodies of written text and could seed early reading material. They are licensed CC BY-ND 4.0, which allows redistribution but not altered text, so adapting them would need permission. Everyday-subject primers, an orthography guide and teacher training would all have to be developed. I did not find evidence of existing vernacular schooling.

\end{document}